% !TEX root = ../../main.tex

\section{系统功能模块设计}

根据第三章的功能性需求分析，HearBridge 系统可以划分为移动端功能模块、后台管理端功能模块、服务端业务模块和手势识别服务模块。
功能模块图主要用于说明系统功能结构和各模块职责，与系统总体架构图不同，功能模块图重点描述系统具备哪些功能，而不是各端之间的通信关系。

系统功能模块如图~\ref{fig:chapter04-function-modules} 所示。

\WhuImageFigure[0.95\textwidth][0.62\textheight]{figures/chapter04/function-modules.png}{系统功能模块图}{fig:chapter04-function-modules}

\subsection{移动端功能模块}

移动端功能模块面向普通用户，主要包括用户认证、手语资源浏览、动作示范展示、实时手势训练、句子视频识别和个人中心等子模块。

用户认证模块负责普通用户登录、注册和退出登录等操作，为后续访问系统功能提供身份状态支持。
手语资源浏览模块负责展示手语分类和手语资源列表，使用户能够选择需要学习或训练的手语动作。
动作示范展示模块负责展示手语资源对应的动作示范内容，帮助用户理解标准动作。
实时手势训练模块负责引导用户对准摄像头完成手势动作，并展示系统返回的识别结果。
句子视频识别模块负责支持用户选择本地手语视频，提交识别并查看原始识别结果和语义修正结果。
个人中心模块负责用户资料查看、信息修改、密码修改和退出登录等操作。

\subsection{后台管理端功能模块}

后台管理端功能模块面向系统管理员，主要包括管理员认证、手势分类管理、手势资源管理、样本管理、模型训练管理、模型版本管理和系统概览等子模块。

管理员认证模块负责后台管理员登录和退出登录，保证后台功能仅允许授权用户访问。
手势分类管理模块用于维护手势资源的分类信息，便于移动端按类别组织训练内容。
手势资源管理模块用于维护具体手势资源，包括资源名称、所属分类、封面或动作示范资源等内容。
样本管理模块用于查看和维护训练样本，为模型训练和样本质量检查提供支撑。
模型训练管理模块用于触发或查看模型训练相关任务。
模型版本管理模块用于维护模型版本信息和发布状态，使识别服务能够加载当前发布模型。
系统概览模块用于展示资源、样本和模型等系统维护状态，帮助管理员快速了解系统整体情况。

\subsection{服务端业务模块}

服务端业务模块是系统的数据管理和业务编排中心，主要包括认证与权限管理、资源管理、样本管理、模型版本管理、文件资源管理、识别请求编排和语义修正调用等子模块。

认证与权限管理模块负责普通用户和管理员的登录状态校验。
资源管理模块负责手势分类和手势资源的数据维护，为移动端展示和管理端操作提供统一接口。
样本管理模块负责保存和查询样本相关信息，为模型训练提供数据来源。
模型版本管理模块负责维护模型版本信息、发布状态和当前可用模型。
文件资源管理模块负责对接对象存储服务，处理头像、封面、示范资源、样本视频和模型文件等非结构化资源。
识别请求编排模块负责接收移动端或管理端的识别和训练相关请求，并按需调用 Python 手势识别服务。
语义修正调用模块负责将句子视频识别结果提交给 DeepSeek，并组织返回给移动端的语义修正结果。

\subsection{手势识别服务模块}

手势识别服务模块主要由 Python 服务承担，负责提供与手势识别和模型处理相关的能力。
该模块主要包括实时手势识别、句子视频识别、样本处理和模型训练与加载等子模块。

实时手势识别模块用于接收移动端相机帧并返回当前识别结果，支撑用户实时训练场景。
句子视频识别模块用于对用户提交的视频进行识别分析，生成词级识别结果。
样本处理模块用于处理采集或导入的手势样本，为后续模型训练提供数据基础。
模型训练与加载模块用于训练识别模型、加载当前发布模型，并向业务系统提供识别能力。
第四章仅对识别服务模块的职责进行设计说明，具体模型结构、特征构造和实验过程将在后续系统实现与测试章节中展开。
